\chapter{Obesity}
\index{Obesity}

\section*{Definition and Overview}
Obesity is a chronic, relapsing, multifactorial disease characterised by excessive or abnormal adiposity that impairs health. It is defined by body mass index (BMI) $\ge$30 kg/m$^2$ (Class I 30--34.9, Class II 35--39.9, Class III $\ge$40), with lower thresholds for Asian populations ($\ge$27.5 kg/m$^2$) and recognition that BMI alone is insufficient. Central adiposity (waist circumference $\ge$94 cm men, $\ge$80 cm women for Europids; $\ge$90/80 cm for Asian) better reflects metabolic risk. Obesity is a disease in its own right and a driver of >200 complications including type 2 diabetes, cardiovascular disease, osteoarthritis, obstructive sleep apnoea, and 13 cancers. It requires long-term, evidence-based management, not simplistic ``eat less, move more'' advice.

\section*{Epidemiology}
In many countries, 67\% of adults (12.5 million) are overweight or obese (BMI $\ge$25), and 31\% (5.8 million) have obesity (BMI $\ge$30). Severe obesity (Class II/III) affects ~10\% and is rising fastest. Childhood obesity affects 25\% of children (2--17 years). Indigenous people have 1.5x higher obesity prevalence. Global obesity has nearly tripled since 1975. Obesity reduces life expectancy by 5--10 years (Class III) and costs the local economy >\$11.8 billion annually in direct healthcare and lost productivity. The COVID-19 pandemic highlighted obesity as an independent risk factor for severe outcomes.

\section*{Aetiology and Pathophysiology}
\begin{itemize}
    \item \textbf{Energy homeostasis dysregulation:} defended body weight set-point elevated; hypothalamic centres (ARC, PVN, LH) integrate peripheral signals (leptin, insulin, ghrelin, GLP-1, PYY, CCK) and regulate intake/expenditure; obesity involves leptin resistance, impaired satiety signalling, and hedonic overeating
    \item \textbf{Genetics:} heritability 40--70\%; monogenic (LEP, LEPR, MC4R, POMC, PCSK1) rare; polygenic (GWAS: >1,000 loci, FTO strongest); gene-environment interaction
    \item \textbf{Environment:} obesogenic -- ultra-processed foods, large portions, sugary drinks, food marketing, built environment discouraging activity, screen time, sleep disruption, endocrine disruptors, antibiotic exposure, socioeconomic factors
    \item \textbf{Adipose tissue dysfunction:} hypertrophy $\to$ hypoxia, fibrosis, inflammation (crown-like structures), macrophage infiltration $\to$ systemic inflammation, insulin resistance, ectopic fat (liver, muscle, pancreas, heart)
    \item \textbf{Gut microbiome:} reduced diversity, altered Firmicutes/Bacteroidetes; increased energy harvest, LPS translocation, inflammation
    \item \textbf{Medications:} antipsychotics, antidepressants, corticosteroids, insulin, sulfonylureas, beta-blockers, antihistamines, gabapentinoids
    \item \textbf{Endocrine/secondary:} hypothyroidism, Cushing, PCOS, hypothalamic injury, growth hormone deficiency -- rare but treatable
\end{itemize}

\section*{Clinical Features}
\begin{itemize}
    \item \textbf{Anthropometric:} BMI, waist circumference, waist-to-height ratio (>0.5 indicates risk), neck circumference
    \item \textbf{Metabolic:} insulin resistance, prediabetes, type 2 diabetes, dyslipidaemia (high TG, low HDL, sdLDL), hypertension, hyperuricaemia
    \item \textbf{Mechanical:} osteoarthritis (knee, hip, spine), obstructive sleep apnoea, GERD, urinary incontinence, hernias, plantar fasciitis, venous insufficiency, lymphedema
    \item \textbf{Cardiovascular:} HFpEF, atrial fibrillation, coronary artery disease, stroke, venous thromboembolism
    \item \textbf{Hepatic:} NAFLD/MAFLD (70--90\% with obesity) $\to$ NASH, cirrhosis
    \item \textbf{Reproductive:} PCOS, infertility, pregnancy complications (GDM, pre-eclampsia, macrosomia), erectile dysfunction, hypogonadism (men)
    \item \textbf{Malignancy:} 13 cancers (oesophageal adenocarcinoma, gastric cardia, colorectal, liver, gallbladder, pancreatic, postmenopausal breast, endometrial, ovarian, renal, meningioma, thyroid, multiple myeloma)
    \item \textbf{Psychosocial:} depression, anxiety, binge eating disorder, weight stigma, discrimination, reduced quality of life, healthcare avoidance
    \item \textbf{Dermatological:} acanthosis nigricans, skin tags, intertrigo, striae, hidradenitis suppurativa
\end{itemize}

\section*{Differential Diagnosis}
\begin{itemize}
    \item \textbf{Secondary obesity:} screen if rapid onset, childhood onset with developmental delay, short stature, dysmorphism, or endocrine signs
    \item \textbf{Lipodystrophy:} generalised/partial fat loss with severe metabolic syndrome; genetic or acquired (HIV, autoimmune)
    \item \textbf{Monogenic obesity syndromes:} Prader-Willi (hypotonia, hyperphagia, hypogonadism), Bardet-Biedl (retinitis pigmentosa, polydactyly, renal), Alström, POMC/LEP/LEPR/MC4R deficiency
    \item \textbf{Medication-induced:} review all medications; switch if possible
    \item \textbf{Edema states:} heart failure, nephrotic syndrome, cirrhosis, lymphedema -- distinguish from adipose tissue
    \item \textbf{Sarcopenic obesity:} high adiposity + low muscle mass/function; worse outcomes; assess by DEXA, grip strength, gait speed
\end{itemize}

\section*{When to Seek Medical Care}
All adults should have BMI and waist circumference measured at least annually. Weight management discussion indicated for BMI $\ge$25 (or $\ge$23 Asian) with comorbidity, or BMI $\ge$30 regardless. Consider referral to weight management service for BMI $\ge$35, or BMI $\ge$30 with complications not controlled in primary care.

\textbf{Urgent assessment for:}
\begin{itemize}
    \item Obesity hypoventilation syndrome (daytime hypercapnia, morning headache, peripheral oedema)
    \item Severe OSA with cardiovascular complications
    \item Acute complications: VTE, acute coronary syndrome, pancreatitis (hypertriglyceridaemia)
    \item Pre-bariatric surgery evaluation
\end{itemize}

\section*{Diagnosis and Investigations}
\begin{itemize}
    \item \textbf{Anthropometry:} height, weight, BMI, waist circumference, hip circumference, neck circumference
    \item \textbf{Body composition:} DEXA (gold standard for fat mass, lean mass, visceral fat), bioimpedance (accessible, less accurate)
    \item \textbf{Metabolic screening:} fasting glucose/HbA1c, lipid profile, liver function (ALT, GGT), renal function, urate, TSH, 25-OH vitamin D, iron studies
    \item \textbf{Cardiovascular risk:} local absolute CVD risk calculator; consider coronary calcium score if intermediate risk
    \item \textbf{Sleep study:} if symptoms of OSA (snoring, witnessed apnoeas, sleepiness)
    \item \textbf{Psychosocial:} PHQ-9, GAD-7, Binge Eating Scale, weight-related quality of life (IWQOL-Lite)
    \item \textbf{Obesity staging:} Edmonton Obesity Staging System (EOSS 0--4) -- incorporates severity of complications and functional impairment; better predicts mortality than BMI alone
    \item \textbf{Genetic testing:} if early-onset severe obesity with hyperphagia, developmental delay, or syndromic features
\end{itemize}

\section*{Management}
\begin{itemize}
    \item \textbf{5As framework (RACGP):} Ask, Assess, Advise, Agree, Assist/Arrange
    \item \textbf{Lifestyle intervention (foundation, all patients):}
    \begin{itemize}
        \item Diet: Mediterranean, DASH, low-carb, very low energy diet (VLED 800 kcal/d, medically supervised), meal replacements; individualised, culturally appropriate, dietitian-led
        \item Physical activity: 150--300 min/week moderate aerobic + 2 resistance sessions; reduce sedentary time; exercise physiologist
        \item Behavioural: self-monitoring, goal-setting (SMART), stimulus control, problem-solving, cognitive restructuring, relapse prevention; psychologist/CBT
        \item Sleep: 7--9 h/night; treat OSA
        \item Stress management: mindfulness, ACT, CBT
    \end{itemize}
    \item \textbf{Pharmacotherapy (BMI $\ge$30 or $\ge$27 with comorbidity, adjunct to lifestyle):}
    \begin{itemize}
        \item GLP-1 receptor agonists: semaglutide 2.4 mg weekly (STEP trials: -15--17\% weight), liraglutide 3 mg daily (-6--8\%); covered in some public systems for T2D; semaglutide 2.4 mg public medicines subsidy for obesity (BMI $\ge$30 or $\ge$27 with comorbidity) from 2024
        \item Dual GIP/GLP-1: tirzepatide 15 mg weekly (SURMOUNT: -20--25\%); public medicines subsidy for T2D; obesity listing anticipated
        \item Naltrexone/bupropion ER: -5--8\%; contraindicated: uncontrolled HTN, seizure disorder, opioid use
        \item Orlistat: -3--5\%; GI side effects; OTC
        \item Phentermine: short-term (12 weeks) only; sympathomimetic; public medicines subsidy authority required
        \item Setmelanotide: for monogenic obesity (POMC, LEPR, BBS)
    \end{itemize}
    \item \textbf{Bariatric surgery (BMI $\ge$40 or $\ge$35 with comorbidity; $\ge$30 with T2D for metabolic surgery):}
    \begin{itemize}
        \item Sleeve gastrectomy (most common): -25--30\% TWL; low complication rate; irreversible
        \item Roux-en-Y gastric bypass: -30--35\% TWL; best for T2D remission, reflux; higher nutritional risk
        \item One-anastomosis gastric bypass (OAGB): similar to RYGB, simpler
        \item Adjustable gastric band: largely abandoned (high reoperation, lower efficacy)
        \item Outcomes: T2D remission 60--80\% (RYGB), 40--60\% (sleeve); CV mortality reduction; durability requires lifelong follow-up
    \end{itemize}
    \item \textbf{Multidisciplinary team:} GP, endocrinologist/obesity physician, dietitian, exercise physiologist, psychologist, bariatric surgeon, nurse coordinator
    \item \textbf{Weight maintenance:} lifelong vigilance; continued pharmacotherapy/surgery follow-up; weight regain is biological, not failure of willpower
\end{itemize}

\section*{Complications}
\begin{itemize}
    \item \textbf{Cardiometabolic:} type 2 diabetes, CVD, stroke, HFpEF, atrial fibrillation, hypertension, dyslipidaemia, NAFLD/NASH
    \item \textbf{Mechanical:} osteoarthritis (knee replacement 20x risk Class III), OSA, GERD, hernias, incontinence, VTE
    \item \textbf{Cancer:} 13 obesity-related cancers; 4--5\% of all cancers attributable
    \item \textbf{Reproductive:} infertility, PCOS, pregnancy complications, male hypogonadism
    \item \textbf{Psychosocial:} depression, anxiety, eating disorders, weight stigma (in healthcare, employment, education), reduced QoL
    \item \textbf{Infection:} impaired immune response, worse COVID-19, influenza outcomes
    \item \textbf{Perioperative:} difficult airway, wound complications, VTE, cardiac events
    \item \textbf{Mortality:} J-shaped curve; lowest at BMI 22--25; Class III obesity: 8--10 years life lost
\end{itemize}

\section*{Prognosis and Outcomes}
Obesity is a chronic disease requiring lifelong management. Realistic goals: 5--10\% weight loss yields significant health benefits (diabetes prevention, BP reduction, lipid improvement, OSA improvement, QoL). Modern pharmacotherapy (GLP-1 RAs, GIP/GLP-1) achieves 15--25\% weight loss, approaching surgical outcomes for some. Bariatric surgery provides most durable weight loss and comorbidity remission but requires lifelong nutritional surveillance. Weight regain is the rule without ongoing treatment; biological drivers (hormonal adaptation, reduced energy expenditure) must be addressed. Stigma reduction and patient-centred care improve engagement and outcomes.

\section*{Prevention and Self-Care}
\begin{itemize}
    \item Prevent weight gain from young adulthood (average 0.5 kg/year); small sustained changes easier than large losses
    \item Breastfeeding reduces childhood obesity risk
    \item Family-based approach: healthy home food environment, active transport, screen limits, shared meals
    \item Follow local Dietary Guidelines; limit ultra-processed foods, sugary drinks, alcohol
    \item Regular physical activity: find enjoyable movement; incidental activity counts
    \item Prioritise sleep (7--9 h); manage stress
    \item Monitor weight (weekly) and waist; act early on small gains
    \item Seek professional help early; avoid fad diets, detoxes, unproven supplements
    \item If on pharmacotherapy: continue long-term; stopping leads to regain
    \item Post-bariatric surgery: lifelong micronutrient supplementation (multivitamin, B12, iron, calcium, D, folate), annual bloods, protein intake 60--80 g/d
    \item Advocate for weight-inclusive healthcare; report stigma
    \item Support organisations: Obesity some countries, Obesity Collective, ANZOS, Weight Issues Network
\end{itemize}

\section*{Sources and Further Reading}
The following reputable sources provide further information for healthcare students and patients seeking reliable, current advice about obesity.

\begin{itemize}
    \item Obesity Australia. \textit{Evidence-based policy and clinical resources.} https://www.obesityaustralia.org/
    \item Australian and New Zealand Obesity Society (ANZOS). \textit{Clinical guidelines and position statements.} https://www.anzos.com/
    \item RACGP. \textit{Obesity prevention and management (SNAP guide).} https://www.racgp.org.au/
    \item Healthdirect Australia. \textit{Obesity.} https://www.healthdirect.gov.au/obesity
    \item World Obesity Federation. \textit{Global obesity data and resources.} https://www.worldobesity.org/
    \item National Health and Medical Research Council (NHMRC). \textit{Clinical practice guidelines for the management of overweight and obesity.} https://www.nhmrc.gov.au/
    \item Diabetes Australia. \textit{Obesity and type 2 diabetes.} https://www.diabetesaustralia.com.au/
    \item National Heart Foundation. \textit{Obesity and cardiovascular risk.} https://www.heartfoundation.org.au/
    \item eviQ. \textit{Bariatric surgery perioperative care.} https://www.eviq.org.au/
\end{itemize}